% =============================================================================
% Lampiran
%
% PPKI: Setiap lampiran di halaman baru. Caption lampiran menggunakan
% format "Lampiran N  Judul" dengan hanging indent. Counter lampiran
% terpisah dari counter chapter dan figure/table.
% =============================================================================

\chapter*{LAMPIRAN}
\addcontentsline{toc}{chapter}{LAMPIRAN}

% -----------------------------------------------------------------------------
% Lampiran 1
% -----------------------------------------------------------------------------
\captionlampiran{Rata-rata dan simpangan baku beberapa sifat fisik dan
  kimia tanah dari 78 contoh tanah di Kebun Percobaan Ciheuleut}
\label{lamp:tanah}

\vspace{4pt}

\begin{center}
\begin{tabular}{lcc}
  \toprule
  Sifat
    & Rata-rata
    & Simpangan baku \\
  \midrule
  Pasir (\%)                        & 47.66 & 23.81 \\
  Lempung (\%)                      & 21.80 & 11.94 \\
  Liat (\%)                         & 30.72 & 18.09 \\
  C-organik (\%)                    & 0.61  & 0.57  \\
  Rapatan isi (mg m$^{-3}$)         & 1.43  & 0.16  \\
  KTK (mek 100 g$^{-1}$ tanah)$^{\mathrm{a}}$
                                     & 18.08 & 17.09 \\
  KAT pada KL (g g$^{-1}$)          & 23.62 & 10.80 \\
  KAT pada TLP (g g$^{-1}$)         & 11.11 & 9.05  \\
  \bottomrule
\end{tabular}
\end{center}

\par\vspace{3pt}\noindent{\footnotesize
  $^{\mathrm{a}}$Banyaknya 70 contoh tanah; KTK: kapasitas tukar kation,
  KAT: kadar air tanah, KL: kapasitas lapang, TLP: titik layu permanen.\par}

\vspace{1.5cm}

% -----------------------------------------------------------------------------
% Lampiran 2
% -----------------------------------------------------------------------------
\captionlampiran{Umur, indeks luas daun, dan hasil biji kering jagung
  yang ditanam pada lima ketinggian tempat}
\label{lamp:jagung}

\vspace{4pt}

\begin{center}
\begin{tabular}{cccc}
  \toprule
  Ketinggian (m dpl) & Umur (hari) & Indeks luas daun & Hasil (ton ha$^{-1}$) \\
  \midrule
  856 & 115 & 3.10 & 5.69 \\
  605 & 106 & 3.09 & 5.43 \\
  400 & 100 & 2.47 & 4.80 \\
  210 &  93 & 2.46 & 4.25 \\
   10 &  88 & 2.12 & 4.03 \\
  \bottomrule
\end{tabular}
\end{center}
